\section{Communication as a natural indirect effect}
\label{sec:nie}

The three diagnostics in common use fail on the dyads of \S\ref{sec:assay}, and
they fail for a shared structural reason: none of them asks the question that
defines communication. Positive signalling looks only at the sender, and is
satisfied by a sender talking to nobody. Causal influence and positive listening
look only at the receiver, and are satisfied by a receiver attending to noise.
Ablation asks what happens when the message is removed, which is a different
question from what the message conveys---and removal takes the receiver outside
the support of its own training distribution, which is precisely why it fires for
a channel that carries nothing.

What communication requires is that the receiver's payoff depend on the sender's
private state \emph{through} the message. That is a mediation question, and it
has a standard causal answer \citep{pearl2016causal,pearl2009causality}.

\paragraph{Definition.} Let $S$ be the sender's private state, $M$ the message it
emits, and $u_R$ the receiver's return. $S$ influences $u_R$ along two kinds of
path: direct ones, since the sender also moves, occupies space and consumes
resources, and the mediated path through $M$. Write $M(s)$ for the message the
sender's policy would emit in state $s$, and $u_R(s, m)$ for the receiver's return
when the world is in state $s$ and the transmitted message is $m$. The
\emph{natural indirect effect} of the message is
\begin{equation}
\mathrm{NIE} \;=\; \mathbb{E}_{s,\,s' \sim p(S)}\big[\, u_R\big(s,\, M(s')\big)
  \;-\; u_R\big(s,\, M(s)\big) \,\big],
\label{eq:nie}
\end{equation}
with $s'$ drawn independently of $s$. Operationally: hold the world at $s$, so
every direct pathway is untouched, and transmit the message the sender
\emph{would have} produced in an independently drawn world. We estimate it by
running matched pairs of episodes that share initial conditions and the
sensory-noise stream, and differencing within pair.

Two properties do the work.

\begin{proposition}[Insensitivity to uninformative channels]
\label{prop:noise}
If $M \perp S$, then $\mathrm{NIE} = 0$.
\end{proposition}
\begin{proof}
If the message does not depend on the sender's state then $M(s')$ and $M(s)$ are
identically distributed for every $s$, so the two terms of Eq.~\ref{eq:nie} have
equal expectation. \qed
\end{proof}

\begin{proposition}[Insensitivity to unattended channels]
\label{prop:deaf}
If the receiver's policy does not depend on $M$, then $\mathrm{NIE} = 0$.
\end{proposition}
\begin{proof}
Then $u_R(s, m) = u_R(s)$ for all $m$, and the difference vanishes pointwise. \qed
\end{proof}

\begin{proposition}[Deletion is not mediation]
\label{prop:ablation}
The ablation contrast $\mathbb{E}[u_R(s, \varnothing)] - \mathbb{E}[u_R(s, M(s))]$,
where $\varnothing$ denotes a removed or zeroed message, can be non-zero even when
$M \perp S$.
\end{proposition}
\begin{proof}
By counterexample, realised in \S\ref{sec:assay}: a receiver whose policy
conditions on $M$ commits to a target on the basis of a message that carries no
information about $S$. Removing $M$ leaves it with an input outside the support of
$M$, so it does not commit at all, and its return falls by \KillNoise. Yet
$M \perp S$, so by Proposition~\ref{prop:noise} the mediated effect is zero and no
information has been lost. \qed
\end{proof}

Propositions~\ref{prop:noise} and \ref{prop:deaf} say the NIE is zero in exactly
the two ways a system can fail to communicate while satisfying a standard
diagnostic, and Proposition~\ref{prop:ablation} says why the payoff test that
looks closest to ours is not equivalent to it. The distinction is that
$M(s')$ is drawn from the message's own marginal: the receiver always sees a
message it could have seen, and only the correspondence between message and world
is broken.

\paragraph{Estimation.} Eq.~\ref{eq:nie} needs $M(s')$, the message the sender's
policy would have produced in a counterfactual world. In a simulator this is
directly available. In the referential game we resample the referent, recompute
the sender's message from it, and transmit that while the true referent and the
food stay where they were. In the naturalistic foraging world the sender's state
is not a single latent variable, so we use the trajectory-level analogue: the
message actually transmitted is replaced by the pulse train the same policy
produced in an independently drawn arena. Marginal rate, burst structure and
inter-pulse-interval statistics are preserved exactly---these are the same trains
that policy emits---while the correspondence between the train and the receiver's
world is destroyed.

\paragraph{What it does not do.} The NIE is a statement about payoff, so it will
report zero for a channel that is informative and attended but whose use happens
not to matter to either party in the environment tested; it measures the
\emph{realised value} of the correspondence, not its existence. It also inherits
the usual mediation caveat: the decomposition into direct and indirect paths is
well defined only if the message can be intervened on without disturbing the
direct paths, which is exactly what the channel algebra of \S\ref{sec:algebra}
buys us and which is not available in an unmodified simulator. Finally, like any
causal-effect estimate it is bounded by statistical power, so a null must be read
against a measured detection floor---which is what the validation dyads provide.
